\SGASubsection{4.5}{Standard position}
\SGARefTarget{sga3:xxvi:subsection:4:5}

In this subsection we briefly indicate how some of the preceding
results generalize.

\begin{sgastatement}{\SGAInlineRefTarget{sga3:xxvi:proposition:4-5-1}Proposition 4.5.1}
Let \(P_1,P_2\) be two parabolic subgroups of the \(S\)-reductive
group \(G\).  The following conditions are equivalent:
\begin{enumerate}[label=\textup{(\roman*)}]
\item \(P_1\cap P_2\) is smooth.
\item \(P_1\cap P_2\) is a subgroup of type \textup{(R)}, or of type
\textup{(RC)}, of \(G\).
\item Fpqc-locally, \(P_1\cap P_2\) contains a maximal torus of \(G\).
\item Zariski-locally, \(P_1\cap P_2\) contains a maximal torus of
\(G\).
\end{enumerate}
When \(S\) is semilocal, these conditions are also equivalent to:
\begin{enumerate}[label=\textup{(\roman*)},start=5]
\item \(P_1\cap P_2\) contains a maximal torus of \(G\).
\end{enumerate}
\end{sgastatement}
\SGARefTarget{sga3:xxvi:prop:4:5:1}

\SGAProof
Clearly, \textup{(iv)} implies \textup{(iii)}, and \textup{(v)}
implies \textup{(iv)} when \(S\) is semilocal; condition \textup{(ii)}
implies \textup{(i)} by Exp.\ XXII, \SGARefLink{sga3:xxii:def:5.2.1}{Definition 5.2.1}.  We shall prove
\[
 \textup{(i)}\Rightarrow\textup{(ii)}
 \Rightarrow\textup{(iii)},
\]
and then show that \textup{(iii)} implies \textup{\SGARefLink{sga3:xxii:def:5.2.1:item:i}{(i)}} and
\textup{(iv)}, as well as \textup{(v)} when \(S\) is semilocal.
Put \(K=P_1\cap P_2\).  By \SGARefLink{sga3:xxvi:lemma:4:1:1}{4.1.1} and \textup{[BT65]}, \SGARefLink{sga3:xxvi:subsection:4:5}{4.5}, every
geometric fiber \(K_{\bar s}\) contains a maximal torus
\(T_{\bar s}\) of \(G_{\bar s}\), is connected, and has, relative to
\(T_{\bar s}\), a closed subset of the roots of \(G_{\bar s}\).
Consequently, if \(K\) is smooth over \(S\), then it is a subgroup of
type \textup{(RC)}, hence a fortiori of type \textup{(R)}; see Exp.\
XXII, \SGARefLink{sga3:xxii:def:5.2.1}{5.2.1} and 5.11.1.  Thus \textup{\SGARefLink{sga3:xxii:def:5.2.1:item:i}{(i)}} and \textup{\SGARefLink{sga3:xxii:def:5.2.1:item:ii}{(ii)}} are
equivalent.

If \(K\) is of type \textup{(R)}, it contains a maximal torus of \(G\)
étale-locally, by Exp.\ XXII, \SGARefLink{sga3:xxii:prop:2.2}{2.2} and Exp.\ XIX, \SGARefLink{sga3:xxvi:proposition:6-1}{6.1}.  Hence
\textup{\SGARefLink{sga3:xxii:def:5.2.1:item:ii}{(ii)}} implies \textup{(iii)}.

Assume \textup{(iii)}, and let us prove that \(K\) is smooth.  By
fpqc descent we may suppose that
\[
 G=(G,T,M,R)
\]
is split, where \(T\) is a maximal torus contained in \(P_1\cap P_2\),
and that there are closed subsets \(R_1,R_2\) of \(R\) such that
\[
 P_1=H_{R_1},\qquad P_2=H_{R_2}.
\]
Since \(K\) has connected fibers, Exp.\ XXII, \SGARefLink{sga3:xxii:prop:5.4.5}{5.4.5} gives
\[
 K=H_{R_1\cap R_2};
\]
therefore \(K\) is smooth over \(S\) and of type \textup{(RC)}.

Let \(R_1^s\) be the symmetric part of \(R_1\), and put
\(R_1^a=R_1-R_1^s\); then
\[
 \rad^u(P_1)=H_{R_1^a}.
\]
As observed in the proof of \textup{[BT65]}, \SGARefLink{sga3:xxvi:subsection:4:4}{4.4},
\[
 R'=(R_1^s\cap R_2)\cup R_1^a
\]
is a closed subset of \(R\) such that \(R'\cup(-R')=R\).  Hence, by
Exp.\ XXI, \SGARefLink{sga3:XXI:3.3.6}{3.3.6}, \(R'\) contains a system of positive roots of \(R\).
Moreover, the symmetric part of \(R'\) is \(R_1^s\cap R_2^s\), which
is contained in \(R_1\cap R_2\).

It follows that
\[
 P'=K\cdot\rad^u(P_1)
\]
is a parabolic subgroup of \(G\), and that \(K\) is a subgroup of
type \textup{(RC)} of \(P'\) satisfying
\[
 \rad^u(K)=K\cap\rad^u(P').
\]
By \SGARefLink{sga3:xxvi:cor:2:11}{2.11}, \(K\) therefore has a Levi subgroup \(L\).  Exp.\ XIV, \SGARefLink{sga3:xiv:3.20}{3.20}
and \SGARefLink{sga3:xxvi:paragraph:3-21}{3.21} then show that \(K\) satisfies \textup{(iv)}, and also
\textup{(v)} when \(S\) is semilocal.  This proves \SGARefLink{sga3:xxvi:proposition:4-5-1}{Proposition 4.5.1}.
\SGAQED
\footnote{\emph{Editor's note.} The proof of the equivalence, together
with condition \textup{(v)}, which was used implicitly in \SGARefLink{sga3:xxvi:cor:6:17}{6.17} of the
original edition, was added in the Polo--Gille edition.  The original
\SGARefLink{sga3:xxvi:proposition:4-5-1}{4.5.1} was consequently divided into \SGARefLink{sga3:xxvi:proposition:4-5-1}{Proposition 4.5.1} and \SGARefLink{sga3:xxvi:def:4:5:1:1}{Definition
4.5.1.1}.}

\begin{sgastatement}{Definition 4.5.1.1}
When the preceding conditions hold, \(P_1\) and \(P_2\) are said to be
in \emph{mutual standard position}.  This is the case, for example,
when they are in transverse position, when they are in osculating
position, or when the base is the spectrum of a field.  The notion is
stable under base extension and is fpqc-local.
\end{sgastatement}
\SGARefTarget{sga3:xxvi:def:4:5:1:1}

\SGAInlineRefTarget{sga3:xxvi:paragraph:4-5-2}\paragraph{\SGARefLink{sga3:xxvi:paragraph:4-5-2}{4.5.2}.}
\SGARefTarget{sga3:xxvi:item:4:5:2}
Let \((P,Q)\) and \((P',Q')\) be two pairs of parabolic subgroups of
\(G\) in standard position.  Define the subfunctor \(H\) of \(G\) by
letting \(H(S')\) be the set of \(g\in G(S')\) such that
\[
 \Int(g)P=P',\qquad \Int(g)Q=Q'.
\]
It is a closed subscheme of \(G\), smooth over \(S\) and formally a
principal homogeneous space under \(P\cap Q\).  It follows that the
following conditions are equivalent:
\begin{enumerate}[label=\textup{(\roman*)}]
\item \((P,Q)\) and \((P',Q')\) are conjugate fpqc-locally;
\item \((P,Q)\) and \((P',Q')\) are conjugate étale-locally;
\item \((P,Q)\) and \((P',Q')\) are conjugate on every geometric
fiber.
\end{enumerate}
The pairs \((P,Q)\) and \((P',Q')\) are then said to have the same
\emph{type of mutual position}.  This notion is stable under base
change and is fpqc-local.

\paragraph{4.5.3.}
\SGARefTarget{sga3:xxvi:item:4:5:3}
Let \(\Stand(G)\) be the subfunctor of
\(\Par(G)\times_S\Par(G)\) formed by the pairs in mutual standard
position.  The functor \(\Stand(G)\) is representable.  There is a
finite étale \(S\)-scheme \(\TypeStand\), the \emph{scheme of types of
mutual standard position}, and a smooth morphism of finite
presentation with irreducible geometric fibers, hence in particular
faithfully flat,
\[
 t_2:\Stand(G)\longrightarrow\TypeStand.
\]
This morphism is a quotient of \(\Stand(G)\) by the action of \(G\):
two sections of \(\Stand(G)\), over any \(S'\to S\), have the same
type of mutual position if and only if they have the same image under
\(t_2\).  There is a commutative diagram
\[
\begin{tikzcd}[column sep=huge,row sep=large]
 \Stand(G) \arrow[r,"t_2"] \arrow[d]
   & \TypeStand \arrow[d,"q"]\\
 \Par(G)\times_S\Par(G) \arrow[r,"t\times t"]
   & \Of(\underline{\operatorname{Dyn}}(G))
       \times_S\Of(\underline{\operatorname{Dyn}}(G)).
\end{tikzcd}
\]
The morphism \(q\) may be described by descent as follows.  Let
\((P,Q)\) be a pair of parabolic subgroups of \(G\) in mutual standard
position, and let \(T\) be a maximal torus of \(P\cap Q\).  The
morphism
\[
 \Norm_G(T)\longrightarrow\Stand(G),\qquad
 n\longmapsto(P,\Int(n)Q),
\]
defined set-theoretically as displayed, induces an isomorphism
\[
 W_P(T)\backslash W_G(T)/W_Q(T)
 \xrightarrow{\sim}q^{-1}\bigl(t(P),t(Q)\bigr).
\]
The left-hand side denotes the sheaf of double cosets.  These
assertions are immediate; note in particular that
\[
 t_2^{-1}\bigl(t_2(P,Q)\bigr)\simeq G/(P\cap Q).
\]

\SGAInlineRefTarget{sga3:xxvi:paragraph:4-5-4}\paragraph{\SGARefLink{sga3:xxvi:paragraph:4-5-4}{4.5.4}.}
\SGARefTarget{sga3:xxvi:item:4:5:4}
Fix a parabolic subgroup \(P\) of \(G\), and let \(\Par(G;P)\) be the
functor of parabolic subgroups of \(G\) in standard position relative
to \(P\).  For every
\[
 t\in\Of(\underline{\operatorname{Dyn}}(G))(S),
\]
put
\[
 \Par_t(G;P)=\Par(G;P)\cap\Par_t(G).
\]
Both functors are obtained from \(\Stand(G)\) by fiber products, and
are therefore represented by \(S\)-schemes, smooth and of finite
presentation over \(S\), with nonempty fibers.  There is a canonical
morphism induced by \(t_2\),
\[
 t_P:\Par_t(G;P)\longrightarrow q^{-1}\bigl(t(P),t\bigr),
 \qquad t_P(Q)=t_2(P,Q),
\]
which is smooth and of finite presentation, with irreducible
geometric fibers.  The canonical morphism
\[
 \Par_t(G;P)\longrightarrow\Par_t(G)
\]
is a surjective monomorphism, and may therefore be viewed as a
\emph{cell decomposition} of \(\Par_t(G)\), indexed by the set of
connected components of \(q^{-1}(t(P),t)\).

\SGAInlineRefTarget{sga3:xxvi:paragraph:4-5-5}\paragraph{4.5.5.}
\SGARefTarget{sga3:xxvi:item:4:5:5}
Suppose that the type \(t\) is of the form \(t(Q)\), where \(Q\) is a
parabolic subgroup of \(G\) in standard position relative to \(P\),
and suppose that \(P\cap Q\) contains a maximal torus \(T\).  Then
\[
 \Par_t(G)\simeq G/Q,\qquad
 q^{-1}\bigl(t(P),t\bigr)
 \simeq W_P(T)\backslash W_G(T)/W_Q(T),
\]
and there is a diagram
\[
\begin{tikzcd}[column sep=huge,row sep=large]
 \Par_{t(Q)}(G;P)
   \arrow[r,"f"] \arrow[d,"i"']
   & W_P(T)\backslash W_G(T)/W_Q(T)\\
 \Par_{t(Q)}(G) \arrow[r,"\sim"]
   & G/Q .
\end{tikzcd}
\]
Here \(i\) is a surjective monomorphism, and \(f\) is smooth and of
finite presentation, with irreducible geometric fibers.  Moreover, if
\(Q_1,Q_2\) are two sections of \(\Par_{t(Q)}(G;P)\), over some
\(S'\to S\), that is, two parabolic subgroups of \(G_{S'}\) which are
fpqc-locally conjugate to \(Q\) and are in standard position relative
to \(P_{S'}\), then \(Q_1,Q_2\) are fpqc-locally conjugate by a
section of \(P\) if and only if
\[
 f(Q_1)=f(Q_2).
\]
If \(S\) is the spectrum of an algebraically closed field \(k\), this
gives
\[
 P(k)\backslash G(k)/Q(k)
 \simeq
 W_P(T)(k)\backslash W_G(T)(k)/W_Q(T)(k).
\]

More generally, suppose that the scheme
\(W_P(T)\backslash W_G(T)/W_Q(T)\) is constant, of the form \(E_S\),
as happens for example when \(G\) is split relative to \(T\) and
\(S\) is connected.  The subschemes \(f^{-1}(e)\), for \(e\in E\),
then decompose \(\Par_{t(Q)}(G;P)\) into open and closed subschemes.
They are homogeneous spaces under \(P\), smooth and of finite
presentation over \(S\), with irreducible geometric fibers.

\SGAInlineRefTarget{sga3:xxvi:paragraph:4-5-6}\paragraph{\SGARefLink{sga3:xxvi:paragraph:4-5-6}{4.5.6}.}
\SGARefTarget{sga3:xxvi:item:4:5:6}
Return to the general situation of \SGARefLink{sga3:xxvi:paragraph:4-5-4}{4.5.4}.  The scheme
\(q^{-1}(t(P),t)\) always has two distinguished sections,
corresponding respectively to the types ``transverse position'' and
``osculating position''.  The inverse image of the first section is
a relatively dense open subscheme of \(\Par_t(G;P)\), as seen above;
it is the cell of greatest relative dimension in the decomposition.

The inverse image of the second section is presumably a closed
subscheme of \(\Par_t(G;P)\); it is the cell of least relative
dimension in the decomposition.
